\documentclass{article}
\usepackage{graphicx}
\usepackage[brazilian]{babel}
\usepackage[latin1]{inputenc}
\usepackage{amssymb}
\usepackage{amsmath}
\usepackage{indentfirst}

\usepackage{xcolor}
\definecolor{verde}{rgb}{0,0.5,0}

\usepackage{listings}
\lstset{
  language=c,
  basicstyle=\ttfamily\small, 
  keywordstyle=\color{blue}, 
  stringstyle=\color{verde}, 
  commentstyle=\color{red}, 
  extendedchars=true, 
  showspaces=false, 
  showstringspaces=false, 
  numbers=left,
  numberstyle=\tiny,
  breaklines=true, 
  backgroundcolor=\color{green!10},
  breakautoindent=true, 
  captionpos=b,
  xleftmargin=0pt,
}

\begin{document}

\title{Lista de Exerc�cios 3 - Filas}
\author{PUC Minas - Engenharia de Computa��o \\ Algoritmos e Estrutura de Dados II \\ Prof. Kleber Jacques F. de Souza}
\date{Data de Entrega: 29/04/2014}
\maketitle

%===================================================================================================

\section*{Instru��es}

\begin{itemize}
\item Esta � uma lista de refor�o.
\item A lista � individual e as datas de entrega s�o fixas, devendo o trabalho ser entregue 
no in�cio da aula. O valor desta lista � 5 pontos. Para cada dia de atraso na entrega 
ser� descontado 0,5 ponto do valor da lista. 
\end{itemize}

%===================================================================================================

\section*{Filas}

\begin{enumerate}
\item Implemente uma fila para gerenciar dados do tipo \textbf{Cliente}, onde Cliente � uma ESTRUTURA com Nome e Idade. Implemente uma para a vers�o usando arranjos e outra para a vers�o usando ponteiros.

\item Implemente um procedimento CLONAR para a Fila, que retorna uma Fila exatamente igual � vari�vel sobre a qual o m�todo foi invocado. Implemente um para a vers�o usando arranjos e outro para a vers�o usando ponteiros.

\item Implemente uma fun��o que receba como par�metro duas Filas e retorne uma terceira contendo a concatena��o das duas. Implemente uma para a vers�o usando arranjos e outra para a vers�o usando ponteiros.

\item Escreva um procedimento que inverta a ordem dos elementos na Fila. Implemente uma para a vers�o usando arranjos e outra para a vers�o usando ponteiros. \underline{DICA:} considera a utiliza��o de estruturas auxiliares, como uma pilha.

\item Em muitas situa��es � necess�rio que exista uma prioridade na Fila, como por exemplo em Bancos onde pessoas idosas t�m prioridade no atendimento. Altere a Fila criada no exerc�cio 1 para que a inser��o na Fila seja feita por prioridade de idade, ou seja, se o Cliente que tiver mais de 65 anos deve ter prioridade no atendimento. Implemente uma para a vers�o usando arranjos e outra para a vers�o usando ponteiros.

\noindent \underline{DICA:} N�o � necess�rio criar a Fila toda novamente, apenas um novo procedimento de Enfileira por prioridade.

\end{enumerate}






%===================================================================================================

\end{document}